% Pro C. Rabirio Perduellionis Reo --- headnote
% pro-rabirio-perduellionis, early/mid 63 BC

\textit{Cicero's defence as consul of the elderly senator
Gaius Rabirius, prosecuted by the tribune Titus Labienus on a
revived charge of \textit{perduellio} --- high treason against
the Roman state --- for an alleged part in the killing of the
tribune Lucius Saturninus thirty-seven years earlier, in 100
BC. Saturninus and his confederate Glaucia had seized the
Capitoline; the consul Marius, by decree of the Senate
(\textit{senatus consultum ultimum}), called all loyal
citizens to arms, and Saturninus was killed after surrender,
together with Labienus's own uncle Quintus, who was on the
Capitoline among the rebels. Labienus's prosecution was
revived in 63 BC at the instigation of Caesar, who saw an
opportunity to put the legitimacy of the senatus consultum
ultimum on trial: if the consul Marius's emergency decree had
not warranted the killing of citizens without trial, then the
weapon by which the Senate had defended itself for a generation
was broken --- and would not be available to Cicero in the
months ahead. The procedural form of \textit{perduellio} was
itself archaic: judgment by two officers (\textit{duumviri
perduellionis}), an executioner waiting in the Forum, the
sentence of crucifixion on the Field of Mars in the words
recited from the regal commentaries (``go, lictor, bind his
hands; veil his head; hang him on the unlucky tree''). Cicero
calls these the carmina of Tarquin, the dirges of torture, and
in his consulship has caused them to be set aside; the
prosecution had to be transferred from the duumviri to a
straight prosecution before the people in the comitia
centuriata, with Cicero limited by Labienus to half an hour
of speech.}

\textit{The argument is consular through and through. Cicero
does not deny that arms were taken up against Saturninus; he
embraces it, and demands that everyone who that day rallied to
the consuls' call --- Marius, Catulus, Scaurus, the two
Mucii, Crassus, Antonius, the entire equestrian order, the
princeps senatus Aemilius leaning on his stick, the crippled
Scaevola leaning on a spear, the praetors and the consulars
and ``all the most illustrious men'' --- be acknowledged as
sharing in Rabirius's act. Either the senatus consultum
ultimum justifies all of them, or it justifies none. The
prosecution is unmasked as an attack not on one elderly
senator but on the supreme protection of the commonwealth, the
\textit{auxilium maiestatis atque imperi}, in the year that
Cicero will need to invoke it against Catiline. The closing
peroration begs the citizens not by their verdict to take from
him, and from the commonwealth, ``the hope of liberty, the
hope of safety, the hope of dignity.''}

\textit{The text survives mutilated: there are textual gaps in
sections 31 and 35, and the closing lines are damaged. The
trial itself ended without a verdict --- when the proceedings
in the comitia centuriata were near judgment, the praetor
Metellus Celer is said to have hauled down the red flag that
flew on the Janiculum (the signal for an enemy attack on the
city) and dissolved the assembly under the old law. Rabirius
went free; the senatus consultum ultimum stood; and a few
months later Cicero invoked it against the Catilinarian
conspirators.}
